\documentclass[11pt]{article}
\usepackage[letterpaper,margin=0.9in]{geometry}
\usepackage[T1]{fontenc}
\usepackage{lmodern}
\usepackage{amssymb}
\usepackage{enumitem}
\usepackage[hidelinks]{hyperref}
\usepackage[protrusion=true,expansion=false]{microtype}
\setlength{\parindent}{0pt}
\setlength{\parskip}{0.45em}
\pagestyle{empty}
\setlength{\emergencystretch}{3em}

\begin{document}

\noindent Davood Damircheli\\
Department of Mathematics\\
Louisiana State University\\
Baton Rouge, LA 70803, USA\\
\href{mailto:ddamircheli@lsu.edu}{ddamircheli@lsu.edu}

November 22, 2026

Editor-in-Chief\\
\emph{SIAM Journal on Scientific Computing}\\
Society for Industrial and Applied Mathematics

\textbf{Re: Submission of ``Symbol-based residual certification of reduced and neural surrogates for
tempered-fractional CGMY pricing''}

Dear Editor-in-Chief,

I am pleased to submit the above manuscript, together with its Supplementary Materials, for
consideration in the section \emph{Numerical Algorithms for Scientific Computing} of SISC.

Fast surrogates of parametric PDE solution maps, whether projection-based reduced models or
neural networks, are only as useful as the guarantee that comes with each prediction. The paper
gives such a guarantee for a demanding class of problems: nonlocal, nonsymmetric parabolic
equations whose fractional order varies with the parameters, here the tempered-fractional
pricing equation of the CGMY model. Every prediction, from any surrogate, is certified query by
query against the spline Galerkin discretisation, and the certificate is cheap enough to be used
online. The key observation is that the discretisation is shift-invariant, so that the symbol
calculus of Toeplitz sections replaces each numerically estimated ingredient of certified
reduced-basis methods by a closed-form one. The main contributions are:
\begin{itemize}[leftmargin=1.6em,itemsep=0.15em,topsep=0.2em]
\item \textbf{Exact assembly.} The Galerkin matrices are Toeplitz sections of an explicit
periodised Fourier symbol; summing its alias series to a proven tail gives the exact Galerkin
system to $10^{-8}$, verified against an independent 40-digit computation.
\item \textbf{A closed-form stability constant.} A L\'evy--Khintchine argument gives a rigorous,
parameter-explicit lower bound for the symmetric part, valid for all $Y\in(0,2)\setminus\{1\}$,
so that no eigenvalue computation or successive constraint method is needed.
\item \textbf{An $O(N)$ certificate and an $N$-independent reduced certificate.} A banded comparison
norm built from the mass and stiffness symbols gives a rigorous bound with explicit effectivity,
and empirical interpolation of the one-variable symbol makes the operator affine (about 36 terms)
with an error transfer that involves no unknown constant. Together with a numerically stable
offline/online evaluation, this certifies reduced models at a cost independent of the resolution,
and a three-stage rule escalates to full order only when needed.
\item \textbf{Evidence at a resolution that prices the whole parameter box to one cent.} The
full-order certificate costs $6\%$ of a direct solve; the certified POD--Galerkin model costs
$3\%$ and settles $99.8\%$ of $10^4$ queries without touching the full space. There is no false
accept on any of the 414 queries with a reference solution, and outside the training box the
certificate accepts 193 of the 210 queries on which the reduced model remains within one cent.
For neural reduced models the certificate is a sharp diagnostic: it rejects every inadequate
prediction within milliseconds and identifies the mechanism. The framework also yields a residual
certificate for the American put.
\end{itemize}
The ingredients (symbol-based assembly and stability constants, banded comparison norms, affine
interpolation of a univariate symbol, query-wise certification of learned surrogates) apply to any
shift-invariant discretisation of a nonlocal operator, including fractional diffusion and other
L\'evy-type models beyond finance, which is why I believe the paper suits SISC.

\textbf{Related work by the author.}
(i) A preprint of this manuscript is available as
\href{https://arxiv.org/abs/2611.03463}{arXiv:2611.03463}. The submitted version contains the same
theory and the same numerical results; it differs from the preprint only by editorial revision
(presentation, length and the Supplementary Materials). The preprint has not been submitted to any
other journal.
(ii) The deterministic spline Galerkin discretisation on which the certificate is built is the
subject of a companion paper, ``A B-spline Galerkin method for option pricing under the CGMY
tempered stable process,'' currently under review at the \emph{SIAM Journal on Financial
Mathematics}. That paper develops and analyses the discretisation; the present one takes it as
given and contributes the certification theory, the reduced and neural surrogates and their
evaluation. The overlap is limited to the model definition, recalled in Section~2 for
self-containedness. I would be glad to provide a copy of the companion manuscript to the editor
and referees.

\textbf{Length and supplementary materials.} The main article has 26 pages, above the recommended
20 and within the 26-page limit. The theory and the numerical evidence depend on each other: the
certificate's cost and reliability claims rest on the calibration, benchmark and out-of-distribution
studies. The main article is self-contained: every result it claims is proved there, including the
uniform-stability theorem and the American-put certificate in Appendix~A. The 8-page Supplementary
Materials contain only supporting data (notation, full verification studies, a neural-operator
stress test, training and cost tables) and three extensions that the main article does not use. I would be happy to shorten the article
further following the editor's guidance.

\textbf{Reproducibility.} All results were produced by a single package with a 67-test suite; a clean
clone regenerates every table and figure of the paper from the stored result files, and every number
in the numerical section is checked against those files by a script. The archive is available at
\url{https://github.com/DavoodDamircheli/ertified_neural_glkn_cgmy}.

\textbf{Suggested referees.} The following experts cover the reduced-basis, Toeplitz/fractional and
computational-finance aspects of the paper:
\begin{itemize}[leftmargin=1.6em,itemsep=0.1em,topsep=0.2em]
\item Bernard Haasdonk, University of Stuttgart, Germany: certified reduced-basis methods and
a posteriori error bounds, including option pricing.
\item Stefano Serra-Capizzano, University of Insubria, Italy, and Uppsala University, Sweden: spectral
symbols of Toeplitz and GLT sequences, and preconditioning of fractional diffusion equations.
\item Olena Burkovska, Oak Ridge National Laboratory, USA: model order reduction for nonlocal and
fractional models, and reduced-basis methods for option pricing.
\item Kathrin Glau, Queen Mary University of London, UK: complexity reduction and parametric
option pricing under L\'evy models.
\item Cornelis W. Oosterlee, Utrecht University, the Netherlands: Fourier methods and machine
learning in computational finance.
\end{itemize}
I have no conflict of interest with any of them.

The manuscript is not published, accepted or under consideration elsewhere, and its posting on
arXiv is consistent with SIAM's policy on preprints. Thank you for considering this submission; I look
forward to the referees' comments.

Sincerely,

\vspace{0.6em}
Davood Damircheli\\
Department of Mathematics, Louisiana State University

\end{document}
